%!TEX program = uplatex
\RequirePackage{plautopatch}
\documentclass[uplatex,dvipdfmx]{jlreq}
\usepackage{amsmath,amssymb}

\pagestyle{empty}
\begin{document}

\noindent
{\large\bfseries トーリック多様体入門}

\medskip
\hspace*{\fill}{\large\bfseries 石井志保子（東工大・理）}

\bigskip\bigskip

$(\mathbf{C}^*)^n = \mathbf{C}^* \times \mathbf{C}^* \times \cdots \times \mathbf{C}^*$（$n$ 個の積）を $n$ 次元トーラスと呼ぶ。
これは積
\[
  (t_1, t_2, \ldots, t_n)(s_1, s_2, \ldots, s_n)
\]
により可換群になる。
いま代数多様体 $X$ にトーラス $(\mathbf{C}^*)^n$ が作用しているとする。
特に $X$ が正規で $(\mathbf{C}^*)^n$ と同型な開軌跡をもつとき、$X$ を $n$ 次元トーリック多様体と呼ぶ。

このトーリック多様体はとても把握しやすい多様体なので、
代数多様体の良いモデルとして、また新しい多様体を具体的に構成する場合にしばしば用いられる。
ではなぜ把握しやすいのか、どうやって把握するのかを紹介するのがこの講演の目的である。

講演では $n$ 次元実線形空間の上に扇と呼ばれる有理多面錐の族を考える。
この $n$ 次元実線形空間の上の扇と $n$ 次元トーリック多様体とは一対一に対応することを紹介する。
そしてトーリック多様体の色々な性質は扇の性質で記述されることを見る。
したがってトーリック多様体は扇によって完全に把握されるということになる。

\end{document}